Pemanfaatan \gls{ai} \textit{agent} untuk otomatisasi alur kerja (\gls{workflow}) berkembang pesat, namun sebagian besar masih bergantung pada \gls{llm} berukuran sangat besar yang berjalan di layanan awan berbayar sehingga menimbulkan persoalan biaya, latensi, serta privasi data. \gls{slm} yang dapat dijalankan secara lokal menjadi alternatif yang menarik, tetapi akurasinya pada tugas penalaran dan pemanggilan \textit{tool} masih tertinggal dibanding model besar. Penelitian ini mengusulkan optimasi akurasi \gls{slm} lokal untuk \gls{ai} \textit{agent workflow automation} melalui kombinasi dua pendekatan: (1) \gls{finetuning} hemat parameter dengan \gls{lora} untuk adaptasi domain, dan (2) \gls{rag} yang diperkuat mekanisme \gls{reranking} agar konteks paling relevan diberikan kepada model yang memiliki keterbatasan panjang konteks. Penelitian akan mengevaluasi kontribusi tiap komponen dan kombinasinya melalui studi ablasi terhadap metrik akurasi tugas, kualitas retrieval, serta efisiensi komputasi. Kontribusi yang diharapkan adalah rancangan \textit{pipeline} \gls{slm} lokal yang lebih akurat namun tetap efisien untuk diterapkan pada otomatisasi alur kerja berbasis agen.

\vspace{0.5\baselineskip}
\noindent\textbf{Kata Kunci}: \textit{small language model}, \textit{LoRA}, \textit{retrieval-augmented generation}, \textit{re-ranking}, \textit{AI agent}
